\section{Ordered limits and Euclidean properties}
\label{sec:ex-limits}

The lattice spacing is sent to zero at fixed volume before the volume is sent to
infinity. A subsequence of pairs $(N_k,p_k)$ tending to infinity in both
coordinates need not give either iterated limit. We record the compactness
argument that produces one common pair of subsequences for all finite
correlation signatures, and then explain which Euclidean properties pass to
the resulting table. The construction is conditional on uniform finite-cutoff
estimates stated below; their derivation from the Wilson trajectory is a
separate input.

Fix one Wilson family $\mathcal C$. Write $a_N$ for its mesh, $P_p$ for its
physical linear period, and $\ell_{N,p}$ for the number of sites in that
period. For an $n$-tuple of field labels $v$,
let $U_{N,p,v}$ denote the finite-lattice correlation, namely the
finite expectation of the chosen insertions. We assume it is multilinear in
each label. The factor $a_N^{4n}$ appears when these values are smeared by a
test function. Let $\mathrm{Cyc}_p$
be the open set of configurations for which no two distinct insertions have
coordinate differences equal to an integer multiple of the physical period
$P_p$ in any coordinate. Write $\mathrm{NoHit}_p$ for configurations in which each distinct pair has
at least one coordinate difference that is not an integer multiple of $P_p$.
This excludes torus-coincident pairs but permits coordinatewise ties. Let $\mathrm{Gen}$ be the set for which every coordinate
projection separates distinct insertions. Let $\mathrm{Conf}$ be the set of
configurations with no coincident points. Thus $\mathrm{Gen}\subset\mathrm{Conf}$;
$\mathrm{Gen}$ also excludes noncoincident coordinate-projection ties, while
$\mathrm{Conf}$ permits them. The table $S$ below is the limiting Euclidean correlation table.

Here are the two uniform estimates used for extraction. In each statement
``eventually'' has its ordinary meaning along the indicated strictly
increasing sequence.

\begin{definition}[Uniform bounds and moduli]
\label{def:ex-limits-estimates}
Fix an ordered pair $f=(f_{\rm mesh},f_{\rm per})$ of strictly increasing
sequences.

\begin{enumerate}
\item A local mesh modulus means that for every fixed period $p$, arity $n$,
  label tuple $v$, compact $K\subset\mathrm{Cyc}_p$, and $\varepsilon>0$,
  there are $\delta>0$ and $j_0$ such that, for $j\ge j_0$ and lattice
  configurations $x,y\in K$ at mesh $a_{f_{\rm mesh}(j)}$,
  \[
  |U_{f_{\rm mesh}(j),p,v}(x)-U_{f_{\rm mesh}(j),p,v}(y)|<\varepsilon
  \quad\hbox{whenever }|x-y|<\delta.
  \]
  The modulus may depend on $p,n,v,K,\varepsilon$.
\item An expanding-volume modulus means that for every $n,v$, compact
  $K\subset\mathrm{Gen}$, and $\varepsilon>0$, there are $\delta>0$ and
  thresholds $k_0$ and $j_0(k)$ such that for $k\ge k_0$, $j\ge j_0(k)$,
  and lattice configurations $x,y\in K$ at mesh $a_{f_{\rm mesh}(j)}$,
  \[
  |U_{f_{\rm mesh}(j),f_{\rm per}(k),v}(x)
       -U_{f_{\rm mesh}(j),f_{\rm per}(k),v}(y)|<\varepsilon
       \quad\hbox{whenever }|x-y|<\delta.
  \]
  The mesh threshold may depend on the fixed volume index $k$.
\item A compact norm bound means both: for every fixed $p,n,v$ and compact
  $K\subset\mathrm{NoHit}_p$, the lattice values $|U_{f_{\rm mesh}(j),p,v}(x)|$
  are bounded by one constant for all $x\in K$ and all sufficiently large
  $j$; and for every $n,v$ and compact $K\subset\mathrm{Conf}$ they are
  bounded by one constant for all sufficiently large pairs in the ordered
  sense $\forall k\ge k_0\;\forall j\ge j_0(k)$.
\end{enumerate}

In the ordered-volume clause, the outer threshold is the period threshold and
the inner mesh threshold may depend on that period. In particular these
assumptions do not assert a single diagonal limit or a modulus uniform over
all periods.
\end{definition}

\begin{proposition}[Common ordered extraction]
\label{prop:ex-limits-extraction}
Fix a countable algebraic basis for the local field-label spaces. Assume the
local mesh modulus, expanding-volume modulus, and compact norm bound of
Definition~\ref{def:ex-limits-estimates} for every finite tuple from this
basis, for the same family $\mathcal C$ and initial order $f$. Then there are
strictly increasing subsequences $t,u$ of the mesh and volume indices. Put
$\rho_k=f_{\rm per}(u(k))$. There are intermediate continuous functions $J_{p,v}$ on each $\mathrm{Cyc}_p$
and a table $S$ continuous on $\mathrm{Gen}$, with the following properties:

\begin{enumerate}
\item for each fixed $p,n,v$, the address-cell interpolants along the mesh subsequence
  $f_{\rm mesh}\circ t$ converge to $J_{p,v}$ uniformly on every compact subset
  of $\mathrm{Cyc}_p$;
\item for each $n,v$, $J_{\rho_k,v}$ converges to $S$ uniformly on
  every compact subset of $\mathrm{Gen}$ as $k\to\infty$;
\item for each separated test function suitable at fixed $p$, the finite-lattice
  smeared correlations converge first along the mesh subsequence to the
  integral against $J_{p,v}$, and
  then along the volume subsequence to the integral against $S$ (extended by
  zero off $\mathrm{Gen}$, which does not change this integral).
\end{enumerate}

All signature limits use the same $t,u$. The mesh subsequence is chosen
independently of the requested test support. The volume limit is only asserted
along $f_{\rm per}\circ u$, and the mesh limit remains an ordinary limit along
the full mesh subsequence at each eligible fixed period.
\end{proposition}

\begin{proof}
We give the compactness argument, including the common-subsequence step.
For each mesh index $j$ and point $x$, let $Q_j(x)$ be the cell representative
selected by the lattice address map, so that $|Q_j(x)-x|\le c_n a_{f_{\rm
mesh}(j)}$ for a constant $c_n$ depending only on the arity. Define the
position-space function $\widetilde U_j(x)=U_{f_{\rm mesh}(j),p,v}(Q_j(x))$.
This function is defined even when $x$ is not a lattice point. Let $K$ be a
compact subset of an open patch and choose a compact neighborhood $K'$ still
inside that patch. For $x,y\in K$ with $|x-y|<\delta/2$, the representatives
$Q_j(x),Q_j(y)$ belong to $K'$ for all sufficiently large $j$, and
\[
 |Q_j(x)-Q_j(y)|\le |x-y|+2c_n a_{f_{\rm mesh}(j)}<\delta.
\]
The lattice modulus on $K'$ therefore makes $\widetilde U_j$ asymptotically equicontinuous on $K$; the rounded
functions need not themselves be continuous. The compact norm bound on $K'$ makes their values bounded there for all
sufficiently large $j$.

The basis-signature set and periods are countable. Choose a countable dense
subset $D$ of each patch. The values $\widetilde U_j(d)$, over all dense points and
signatures for the fixed period, lie eventually in bounded closed disks in
$\mathbb C$. A diagonal subsequence makes all these scalar sequences converge.
We check that this gives compact-uniform convergence, not just convergence on
$D$. Let $K$ be compact in the patch and $\varepsilon>0$. Choose a compact
neighborhood $K'$ inside the patch. Use the common modulus on $K'$ to choose a
radius $\delta$ for error $\varepsilon/3$. Finitely many balls of radius
$\delta/3$ centered at points of $D\cap K'$ cover $K$. At these finitely many
centers the selected subsequence is Cauchy. For any $x\in K$, choose one such
center $d$; for all large indices $i,i'$, the two modulus errors
$|\widetilde U_i(x)-\widetilde U_i(d)|$ and
$|\widetilde U_{i'}(x)-\widetilde U_{i'}(d)|$ are each less than
$\varepsilon/3$, while the center values differ by less than
$\varepsilon/3$. The triangle inequality proves the subsequence is uniformly
Cauchy on $K$. Completeness of $\mathbb C$ gives a uniform limit on $K$; limits
agree on overlaps, and the same modulus estimate proves the resulting
function is continuous. Exhausting the patch by compact sets gives one
continuous limit there and uniform convergence on every compact subset.
Diagonalize over the countable collection of signatures and a compact
exhaustion of each patch to obtain one mesh subsequence $t$ and all the
intermediate functions $J_{p,v}$.

Now consider the sequence of mesh-limit functions $J_{f_{\rm per}(k),v}$ on
$\mathrm{Gen}$. For a compact $K\subset\mathrm{Gen}$ choose a compact neighborhood
$K'\subset\mathrm{Gen}$ with $K$ in its interior. For all sufficiently large $k$, $K'$ lies inside the fundamental box; since $K'\subset\mathrm{Gen}$, this gives $K'\subset\mathrm{Cyc}_{f_{\rm per}(k)}$. If $x,y\in K$ and
$|x-y|<\delta/2$, their representatives at mesh $a_{f_{\rm mesh}(j)}$ lie
in $K'$ and have mutual distance less than $\delta$ for all sufficiently
large $j$. Apply the expanding-volume modulus at the fixed outer index $k$,
then pass $j\to\infty$ using the already established mesh convergence on
$K'$. This proves asymptotic equicontinuity of the functions $J_{f_{\rm per}(k),v}$ on
$K$. Similarly, the ordered compact norm bound and mesh convergence give a
uniform bound for these functions on $K$ for all sufficiently large $k$.

Choose countable dense subsets in each generic configuration space and
repeat the closed-disk diagonal argument, now with $k$ as the sequence index.
The finite-cover argument proves compact-uniform convergence on every compact
subset of $\mathrm{Gen}$ and gives a continuous limit $S$. Diagonalizing over
the countable signature index makes the same volume subsequence $u$ work for
all signatures. Use arity together with finite tuples from a
countable algebraic basis as the countable signature index. Every arbitrary finite tuple of labels is a finite
linear combination of such basis tuples, and finite lattice moments are
multilinear, so finite linearity extends these limits to arbitrary labels
without any further subsequence. This is a diagonal argument over the countable basis index, not over the
uncountably many arbitrary test functions.

It remains to pass from pointwise limits to smeared moments. Let $\varphi$ be
a separated test, meaning its support $K$ is compact and contained in
$\mathrm{Conf}$, and fix an eligible period $p$. Assume $p$ is eligible for this test, so $K$ lies strictly inside the
comparison cell. Choose a compact neighborhood $K_1$ of $K$ whose closure
stays inside the comparison cell. For all fine meshes, every cell representative of a point
in $K$ lies in $K_1$. The cell boundaries in this bounded region lie in finitely many affine
hyperplanes and have Lebesgue measure zero. The complement of the cyclic patch is a countable union of affine
hyperplanes and therefore has Lebesgue measure zero. At almost every point
in $K$, the cell representatives tend to that point and eventually belong to
the cyclic patch. Extend $J_{p,v}$ by zero off that patch. The mesh
convergence gives almost-everywhere convergence of the cell density to
$\varphi J_{p,v}$.
Since $K$ is compactly contained in the fundamental cell and in $\mathrm{Conf}$,
it lies in the fixed-period no-hit set. Choose a compact neighborhood $K_2$
of $K$ contained in that open set. For all sufficiently fine meshes the rounded samples of points
in $K$ lie in $K_2$. The fixed-period compact norm bound on $K_2$ bounds the finite-lattice
moments for all sufficiently fine meshes. If $\varphi(Q_N(x))\ne0$, the
cell representative $Q_N(x)$ belongs to $K$ and hence to $K_2$. Therefore the
absolute value of the piecewise constant cell density is at most
$D\|\varphi\|_\infty$ on the whole fundamental box and is zero outside
the union of cells meeting the support. The indicator of the bounded
fundamental box times this constant is integrable. The exact cell identity
writes the smeared sum as the integral of this density: each cell has volume
$a_N^{4n}$. The strict cell margin keeps all samples inside the no-wrap
region. At points outside the closed support $K$, the cell density is eventually zero because $Q_N(x)\to x$; outside the fixed box it is identically zero. Dominated convergence proves the fixed-period smeared limit.

For the volume limit, choose compact neighborhoods $K_0\subset K_1\subset
\mathrm{Conf}$ of the support $K$. For sufficiently fine meshes, every cell
representative of a point in $K_0$ lies in $K_1$. The ordered compact norm
bound on $K_1$ gives one constant for these lattice values once the period
index is sufficiently large and then the mesh threshold is met. For all
sufficiently large period indices, $K_1$ lies inside the fundamental cell.
At generic points $x\in K$, the no-wrap equivalence identifies the cyclic
patch with $\mathrm{Gen}$, and the mesh evaluation lemma passes the lattice
bound to $J_{f_{\rm per}(k),v}(x)$. The set $K\setminus\mathrm{Gen}$ is
contained in finitely many coordinate hyperplanes and has Lebesgue measure
zero. Therefore these intermediate functions are bounded by the same constant
almost everywhere on $K$. Extend $S$ by zero off $\mathrm{Gen}$. The
volume-point convergence holds at each $x\in K\cap\mathrm{Gen}$, hence almost
everywhere on $K$. Dominated convergence gives the
second integral limit for every separated test, including supports that meet
noncoincident coordinate-projection ties. Here $\mathrm{Gen}$ supplies a
full-measure set for the integral passage, while the larger set
$\mathrm{Conf}$ contains the entire test support.
\end{proof}

The two compactness stages use distinct moduli: one at each fixed period and
one in the ordered volume limit. The argument preserves the order of limits:
no estimate uniform in an unbounded volume threshold is substituted for the
permitted threshold $j_0(k)$.

\begin{proposition}[Passage of Euclidean identities and positivity]
\label{prop:ex-limits-euclidean}
Let $S$ be a table obtained with the ordered smeared limits of
Proposition~\ref{prop:ex-limits-extraction}. Assume in addition that:

\begin{enumerate}
\item finite-lattice correlations satisfy exact permutation and lattice
  translation identities, signed coordinate symmetries, conjugation, and
  normalization at admissible grid points; labels transform accordingly;
\item the ordered point limits exist on the generic domain and $S$ is
  continuous there; assume also that for each signature it extends to a smooth
  function on $\mathrm{Conf}$ agreeing with $S$ on $\mathrm{Gen}$;
\item for each finite list of positive-time smeared observables, the exact
  finite-lattice pairing identity expresses each reflected product as a finite
  linear combination of separated smeared correlations covered by the
  preceding extraction, with coefficients independent of the lattice indices;
  the same combination defines the limiting smeared pairing. The finite-lattice
  smeared reflection form is nonnegative for all sufficiently late indices in
  the ordered sense.
\end{enumerate}

Then the corresponding permutation, translation, signed coordinate-permutation,
and conjugation identities hold for $S$ on the noncoincident domain, and the
smeared reflection form of $S$ is nonnegative for every finite list of
admissible positive-time test observables. If, in addition, $S$ satisfies the separate cap bound required by the
Euclidean hierarchy interface, these conclusions together with normalization
give its Euclidean positivity and symmetry clauses.
\end{proposition}

\begin{proof}
For permutation and conjugation, fix the signature and a generic configuration
and approximate it by admissible lattice configurations. At each sufficiently
fine mesh the finite-lattice change of variables gives the exact identity;
the point limits along the common ordered filter and uniqueness of limits in
$\mathbb C$ give the identity for $S$. A signed coordinate permutation
preserves the lattice, so the same argument applies after transforming the
grid points and the insertion labels. For translation by an arbitrary vector
$z\in E_4$, choose lattice translations $z_j$ converging to $z$ as the mesh
tends to zero. The finite trajectory's anchoring identity gives equality
between the two grid evaluations separated by $z_j$. The grid points converge
to the desired translated configurations, and continuity of the limiting
correlation gives the translation identity. The generic set is dense in the
noncoincident domain; if an identity has first been proved only there,
continuity of both sides extends it throughout that domain.

For positivity, fix $m$ positive-time smeared observables $A_1,\ldots,A_m$ and
coefficients $c_1,\ldots,c_m\in\mathbb C$. The pairing assumption and the
ordered convergence of each separated entry show that every entry in the
finite reflection matrix converges to its limiting smeared pairing. At each
sufficiently late admissible lattice pair, reflection positivity gives
\[
  Q_{N,p}:=\sum_{i,j=1}^m \overline{c_i}c_j\,
       \langle \theta A_i\,A_j\rangle_{N,p}\ \geq\ 0.
\]
There are only $m^2$ terms. Since every matrix entry has the same ordered
limit, finite-sum continuity gives
\[ Q_{N,p}\longrightarrow\sum_{i,j}\overline{c_i}c_j\langle\theta A_i A_j\rangle_S. \]
This convergence holds along the same ordered filter.
The nonnegative real half-line is closed, so the limit is nonnegative. Since
the list and coefficients were arbitrary, the limiting smeared reflection
matrix is positive semidefinite. Pointwise reflection positivity
for point insertions then follows from the separate bump-approximation argument
in Lemma~\ref{ex-rp-bumps}; that step uses continuity away from coincident
points and is not a consequence of generic point convergence alone.

These limit arguments establish only the clauses just named. The extraction
gives continuity on $\mathrm{Gen}$; passing identities to $\mathrm{Conf}$
requires the separate smoothness hypothesis there. Smoothness on $\mathrm{Conf}$
includes noncoincident coordinate-projection ties but excludes actual
coincident points. Field quotients and Wightman reconstruction do not follow
from a subsequence argument; they are separate hypotheses or results in the
hierarchy and reconstruction units.
\end{proof}

\begin{remark}[Scope of the extraction]
\label{rem:ex-limits-scope}
The local estimates are not a proof that the Wilson family has a convergent
subsequence by themselves. One must first establish the exact norm and oscillation
bounds for that family, uniformly over the compact sets and with the specified
fixed-period or ordered-volume quantifiers. Also, the extracted first surface
is only continuous on generic patches. The later proof of smoothness on the
full noncoincident set (including noncoincident coordinate-projection ties),
followed by the full Euclidean hierarchy, has its own proof obligations.
\end{remark}
